% !TeX program = xelatex
% Compile with XeLaTeX in Overleaf.


\documentclass[a4paper]{article}
\usepackage{xeCJK}
\setCJKmainfont[BoldFont=FandolSong-Bold.otf]{FandolSong-Regular.otf}
\usepackage{needspace}

% \usepackage[default]{fontsetup}

\usepackage{fancyhdr}
\usepackage{extramarks}[=v4] % The homework template uses version 4 marks.
\usepackage{amsmath}
\usepackage{amsthm}
\usepackage{amsfonts}
\usepackage{tikz}
\usepackage[plain]{algorithm}
\usepackage{algpseudocode}
\usepackage{enumerate}
\usepackage{tikz}

\usetikzlibrary{automata,positioning}

%
% Basic Document Settings
%  

\topmargin=-0.2in
\evensidemargin=0in
\oddsidemargin=0in
\textwidth=6.5in
\textheight=9.5in
\headsep=0.25in

\linespread{1.1}

\pagestyle{fancy}
\lhead{}
\chead{\hmwkClass : \hmwkTitle}
\rhead{\firstxmark}
\lfoot{\lastxmark}
\cfoot{\thepage}

\renewcommand\headrulewidth{0.4pt}
\renewcommand\footrulewidth{0.4pt}

\setlength\parindent{0pt}

%
% Create Problem Sections
%

\newcommand{\enterProblemHeader}[1]{
    \nobreak\extramarks{}{Problem \arabic{#1} continued on next page\ldots}\nobreak{}
    \nobreak\extramarks{Problem \arabic{#1} (continued)}{Problem \arabic{#1} continued on next page\ldots}\nobreak{}
}

\newcommand{\exitProblemHeader}[1]{
    \nobreak\extramarks{Problem \arabic{#1} (continued)}{Problem \arabic{#1} continued on next page\ldots}\nobreak{}
    \stepcounter{#1}
    \nobreak\extramarks{Problem \arabic{#1}}{}\nobreak{}
}

\newcommand*\circled[1]{\tikz[baseline=(char.base)]{
		\node[shape=circle,draw,inner sep=2pt] (char) {#1};}}


\setcounter{secnumdepth}{0}
\newcounter{partCounter}
\newcounter{homeworkProblemCounter}
\setcounter{homeworkProblemCounter}{1}
\nobreak\extramarks{Problem \arabic{homeworkProblemCounter}}{}\nobreak{}

%
% Homework Problem Environment
%
% This environment takes an optional argument. When given, it will adjust the
% problem counter. This is useful for when the problems given for your
% assignment aren't sequential. See the last 3 problems of this template for an
% example.
%

\newenvironment{homeworkProblem}[1][-1]{
    \ifnum#1>0
        \setcounter{homeworkProblemCounter}{#1}
    \fi
    \section{Problem \arabic{homeworkProblemCounter}}
    \setcounter{partCounter}{1}
    \enterProblemHeader{homeworkProblemCounter}
}{
    \exitProblemHeader{homeworkProblemCounter}
}

%
% Homework Details
%   - Title
%   - Class
%   - Due date
%   - Name
%   - Student ID

\newcommand{\hmwkTitle}{Homework\ \#07}
\newcommand{\hmwkClass}{Probability \& Statistics for EECS}
\newcommand{\hmwkDueDate}{2026-11-24}
\newcommand{\hmwkAuthorName}{杨涵}
\newcommand{\hmwkAuthorID}{2025531064}


%
% Title Page
%

\author{
	Name: \textbf{\hmwkAuthorName} \\
	Student ID: \hmwkAuthorID}
\date{}

\renewcommand{\part}[1]{\textbf{\large Part \Alph{partCounter}}\stepcounter{partCounter}\\}

%
% Various Helper Commands
%

% Useful for algorithms
\newcommand{\alg}[1]{\textsc{\bfseries \footnotesize #1}}
% For derivatives
\newcommand{\deriv}[1]{\frac{\mathrm{d}}{\mathrm{d}x} (#1)}
% For partial derivatives
\newcommand{\pderiv}[2]{\frac{\partial}{\partial #1} (#2)}
% Integral dx
\newcommand{\dx}{\mathrm{d}x}
% Alias for the Solution section header
\newcommand{\solution}{\par\addvspace{\baselineskip}\Needspace{4\baselineskip}\noindent\textbf{\large Solution}\par\nobreak\smallskip}
% Probability commands: Expectation, Variance, Covariance, Bias
\newcommand{\E}{\mathrm{E}}
\newcommand{\Var}{\mathrm{Var}}
\newcommand{\Cov}{\mathrm{Cov}}
\newcommand{\Bias}{\mathrm{Bias}}

\begin{document}


% \maketitle
% \thispagestyle{empty}
% \pagebreak

\date{
Due on Nov. 24, 2026, 18:00 UTC+8}
\title{SI 140A-01  Probability \& Statistics for EECS, Fall 2026 \\
Homework 7}
\maketitle
Read all the instructions below carefully before you start working on the assignment, and before you make a submission.
\begin{itemize}
    \item You are required to write down all the major steps towards making your conclusions; otherwise you may obtain limited points of the problem.
    \item Write your homework in English; otherwise you will get no points of this homework.
    \item Any form of plagiarism will lead to $0$ point of this homework. 
\end{itemize}
\newpage

\begin{homeworkProblem}[1] \textbf{(30pts)}
Let $U_1,\dots,U_n$ be i.i.d. $\text{Unif}(0,1)$, and $X = \max(U_1,\dots,U_n)$. What is the PDF of $X$? What is $E X$?

\smallskip

Hint: Find the CDF of $X$ first, by translating the event $X \leq x$ into an event involving $U_1,\dots,U_n$

\solution
Let $F_X$ be the CDF of $X$. For $0<x<1$, independence gives
\begin{align*}
F_X(x)
&=P(X\leq x)=P(U_1\leq x,\dots,U_n\leq x)\\
&=\prod_{i=1}^{n}P(U_i\leq x)=x^n.
\end{align*}
Since $X$ takes values in $(0,1)$, its CDF is
\[
F_X(x)=
\begin{cases}
0, & x\leq 0,\\
x^n, & 0<x<1,\\
1, & x\geq 1.
\end{cases}
\]
Differentiating on each interval gives the PDF
\[
f_X(x)=
\begin{cases}
nx^{n-1}, & 0<x<1,\\
0, & \text{otherwise}.
\end{cases}
\]
The expectation is
\begin{align*}
\E[X]
&=\int_0^1 x f_X(x)\,\dx
=\int_0^1 x\cdot nx^{n-1}\,\dx\\
&=n\int_0^1 x^n\,\dx
=n\left[\frac{x^{n+1}}{n+1}\right]_0^1
=\frac{n}{n+1}.
\end{align*}
\end{homeworkProblem}
\newpage

\begin{homeworkProblem}[2] \textbf{(30pts)}
Let $Z \sim \mathcal{N}(0,1)$ and let $S$ be a random sign independent of $Z$, i.e., $S$ is $1$ with probability $1/2$ and $-1$ with probability $1/2$. Show that $SZ \sim \mathcal{N}(0,1)$.

\solution
Let $\varphi$ and $\Phi$ denote the standard normal PDF and CDF. Then
\[
\Phi(x)=P(Z\leq x)=\int_{-\infty}^{x}\varphi(t)\,\mathrm{d}t.
\]
By symmetry and continuity of the standard normal distribution,
\[
\Phi(-x)=P(Z\leq -x)=P(Z\geq x)=1-\Phi(x).
\]

Let $Y=SZ$. Conditioning on $S$ gives
\begin{align*}
F_Y(x)
&=P(SZ\leq x)\\
&=P(SZ\leq x\mid S=1)P(S=1)+P(SZ\leq x\mid S=-1)P(S=-1)\\
&=\frac12\bigl[P(SZ\leq x\mid S=1)+P(SZ\leq x\mid S=-1)\bigr].
\end{align*}
Since $S$ and $Z$ are independent,
\begin{align*}
P(SZ\leq x\mid S=1)
&=P(Z\leq x\mid S=1)=P(Z\leq x)=\Phi(x),\\
P(SZ\leq x\mid S=-1)
&=P(-Z\leq x\mid S=-1)=P(Z\geq -x)\\
&=1-\Phi(-x)=\Phi(x).
\end{align*}
Thus $F_Y(x)=\frac12\bigl[\Phi(x)+\Phi(x)\bigr]=\Phi(x)$ for all $x$,
so $SZ\sim\mathcal{N}(0,1)$.
\end{homeworkProblem}

\newpage
\begin{homeworkProblem}[3] \textbf{(40pts)}
The \emph{Gumbel distribution} is the distribution of $-\log X$ with $X \sim \text{Expo}(1)$.

\begin{enumerate}
    \item Find the CDF of the Gumbel distribution.
    \item Let $X_1, X_2, \dots$ be i.i.d. $\text{Expo}(1)$ and let $M_n = \max(X_1, \dots, X_n)$. Show that $M_n - \log n$ converges in distribution to the Gumbel distribution, i.e., as $n \to \infty$ the CDF of $M_n - \log n$ converges to the Gumbel CDF.
\end{enumerate}

\solution
For $X\sim\text{Expo}(1)$, the PDF is
\[
f_X(t)=
\begin{cases}
e^{-t}, & t>0,\\
0, & t\leq 0.
\end{cases}
\]
For $t>0$, integrating the density gives
\[
F_X(t)=P(X\leq t)=\int_0^t e^{-u}\,\mathrm{d}u
=\left[-e^{-u}\right]_0^t=1-e^{-t}.
\]
For $t\leq 0$, $F_X(t)=0$.

\begin{enumerate}[(i)]
    \item Let $Y=-\log X$. For any $y\in\mathbb{R}$,
    \begin{align*}
    F_Y(y)
    &=P(-\log X\leq y)=P(\log X\geq -y)\\
    &=P(X\geq e^{-y})=e^{-e^{-y}}.
    \end{align*}

    \Needspace{11\baselineskip}
    \item For $t>0$, independence gives
    \[
    P(M_n\leq t)=\prod_{i=1}^{n}P(X_i\leq t)=(1-e^{-t})^n.
    \]
    For each fixed $y\in\mathbb{R}$, $y+\log n>0$ for all sufficiently large $n$.
    For those $n$,
    \begin{align*}
    P(M_n-\log n\leq y)
    &=P(M_n\leq y+\log n)\\
    &=\bigl(1-e^{-(y+\log n)}\bigr)^n
    =\left(1-\frac{e^{-y}}{n}\right)^n.
    \end{align*}
    Taking the limit gives
    \[
    \lim_{n\to\infty}\left(1-\frac{e^{-y}}{n}\right)^n=e^{-e^{-y}}.
    \]
    Thus the CDF of $M_n-\log n$ converges to the Gumbel CDF at every $y$.
\end{enumerate}
\end{homeworkProblem}

\newpage
\begin{homeworkProblem}[4]\textbf{(optional)}
Let $Y \sim \mathcal{N}(\mu, 1)$. Use the fact that $P(|Y - \mu| < 1.96\sigma) \approx 0.95$ to construct a random interval $(a(Y), b(Y))$ (that is, an interval whose endpoints are r.v.s), such that the probability that $\mu$ is in the interval is approximately $0.95$. This interval is called a \emph{confidence interval} for $\mu$; such intervals are often desired in statistics when estimating unknown parameters based on data.

\solution
Since $\sigma^2=1$, we have $\sigma=1$. Rearranging the inequalities gives
\[
|Y-\mu|<1.96
\quad\Longleftrightarrow\quad
-1.96<Y-\mu<1.96
\quad\Longleftrightarrow\quad
Y-1.96<\mu<Y+1.96.
\]
Choose $a(Y)=Y-1.96$ and $b(Y)=Y+1.96$. Then
\[
P\bigl(a(Y)<\mu<b(Y)\bigr)=P(|Y-\mu|<1.96)\approx 0.95.
\]
The resulting confidence interval is $(Y-1.96,Y+1.96)$.
\end{homeworkProblem}
 
 \newpage
 \begin{homeworkProblem}[5]
    Let $F$ be a CDF which is continuous and strictly increasing.
    Let $\mu$ be the mean of the distribution.
    The quantile function, $F^{-1}$, has many applications in statistics and econometrics. 
    Show that the area under the curve of the quantile function from $0$ to $1$ is $\mu$.

\solution
Let $U\sim\text{Unif}(0,1)$ and define $Y=F^{-1}(U)$. Since $F$ is continuous
and strictly increasing,
\[
P(Y\leq y)=P(F^{-1}(U)\leq y)=P(U\leq F(y))=F(y).
\]
Thus $Y$ has CDF $F$ and $\E[Y]=\mu$. Also, using the uniform density $f_U$,
\[
\E[Y]=\E[F^{-1}(U)]
=\int_0^1 F^{-1}(u)f_U(u)\,\mathrm{d}u
=\int_0^1 F^{-1}(u)\,\mathrm{d}u.
\]
Therefore $\int_0^1 F^{-1}(u)\,\mathrm{d}u=\mu$.
\end{homeworkProblem}



\end{document}
